\documentclass[aspectratio=169]{beamer}

\usepackage[UTF8]{ctex}
\usepackage{amsmath,amssymb}
\usepackage{booktabs}
\usepackage{tabularx}
\usepackage{tikz}
\usetikzlibrary{arrows.meta,calc,positioning,fit}
\usepackage[backend=biber,style=authoryear,maxcitenames=2]{biblatex}
\addbibresource{refs.bib}

\definecolor{nslPurple}{HTML}{642860}
\usetheme[
  accent=nslPurple,
  progressbar=foot,
  sectionpage=progressbar,
  subsectionpage=progressbar,
  numbering=fraction,
  numberrail=subsection,
  numberrailcolor=muted,
  block=fill,
  covered=transparent,
  notes=hide,
  toc=aftertitle,
  closing=contact
]{cookie}

\title[双打 OU 对战智能体]{基于强化学习的双打OU对战智能体训练与策略优化研究}
\subtitle{面向高强度竞技环境的自博弈学习、队伍协同建模与强劲AI对手构建}

\author[盖欧卡\quad 固拉多\quad 故勒顿\quad 密勒顿\quad 太乐巴戈斯\quad 烈空坐\textsuperscript{*}]{盖欧卡\textsuperscript{1}\quad 固拉多\textsuperscript{1}\quad 故勒顿\textsuperscript{2}\quad 密勒顿\textsuperscript{2}\quad 太乐巴戈斯\textsuperscript{3}\quad 烈空坐\textsuperscript{*1}}

\institute[丰缘生态研究院 / 帕底亚零区实验室 / 零之大空洞研究所]{%
  \scriptsize
  1. 丰缘地区海陆空环境与传说生态研究院\quad 丰缘地区琉璃市天空柱大道1号\quad 114514\\
  2. 帕底亚第零区零区实验室古今生态与时空演化研究中心\quad 帕底亚地区第零区零区实验室\quad 191981\\
  3. 帕底亚第零区零之大空洞结晶生命研究所\quad 帕底亚地区第零区零之大空洞\quad 191987
}
\date{2026年7月}
\cookieconference{这里放会议名称/硕博论文答辩}
\cookieheaderconference{这里放会议名称/硕博论文答辩}
\cookielogo{assets/文献中心logo.png}
\cookiefootlogo{assets/文献中心logo_白色.png}
\cookieemail{rayquaza@sky-pillar.example}
\cookieclosingtitle{感谢观看}
\cookieclosingtext{欢迎交流与批评指正}

\begin{document}

\maketitle

\section{模板全局结构}
\subsection{封面元数据}

\begin{frame}{封面信息与机构标识}{标题、副标题、作者、单位、会议名和 Logo 的组合}
  \begin{columns}[T,onlytextwidth]
    \column{.48\textwidth}
      \begin{block}{首页展示的信息}
        封面使用顶部会议栏、中心 Logo、标题区、作者区、单位区和日期区。标题区支持副标题；作者和单位可通过 Beamer 的短标题参数控制正文页与结尾页的显示长度。
      \end{block}
    \column{.48\textwidth}
      \begin{cookiecard}[title=元数据命令]
        \small
        \texttt{\textbackslash title[短标题]\{论文题名\}}\\
        \texttt{\textbackslash subtitle\{副标题\}}\\
        \texttt{\textbackslash author[短作者]\{作者列表\}}\\
        \texttt{\textbackslash institute[短单位]\{单位列表\}}\\
        \texttt{\textbackslash cookieconference\{会议名\}}
      \end{cookiecard}
  \end{columns}

  \note{说明封面字段与短标题字段的关系，提醒正式使用时替换虚构作者与单位。}
\end{frame}

\subsection{导航逻辑}

\begin{frame}{顶栏和底栏的显示逻辑}{多级章节驱动结构导航}
  \small
  \begin{columns}[T,onlytextwidth]
    \column{.42\textwidth}
      \begin{nslrule}[结构输入]
        使用 \texttt{\textbackslash section} 和 \texttt{\textbackslash subsection} 组织内容后，模板会自动生成章节页、小节页、正文顶栏和底栏导航。
      \end{nslrule}

      \begin{enumerate}
        \item 顶栏左侧显示会议名称。
        \item 顶栏中间显示当前章节。
        \item 顶栏右侧显示当前小节和小节圆点。
      \end{enumerate}

    \column{.54\textwidth}
      \begin{cookiecard}[title=本页可观察的导航元素]
        \begin{itemize}
          \item \textbf{页眉}: 会议名、章节标题、小节标题和小节进度。
          \item \textbf{页脚}: 文献中心白色 Logo、章节圆点导航和页码。
          \item \textbf{数字轨}: 帧标题左侧以“章节.小节”提示当前位置。
        \end{itemize}
      \end{cookiecard}
      \vspace{.25cm}
      \begin{exampleblock}{演示目的}
        后续章节包含多个 section/subsection，用来观察底栏章节圆点如何随章节切换而变化。
      \end{exampleblock}
  \end{columns}
\end{frame}

\section{内容组件}
\subsection{Block 与列表}

\begin{frame}{Block 族与列表层级}{普通、警示、示例三类块}
  \begin{columns}[T,onlytextwidth]
    \column{.31\textwidth}
      \begin{block}{普通 Block}
        用于承载一般性说明、定义前置条件和页面中的稳定信息。
      \end{block}

    \column{.31\textwidth}
      \begin{alertblock}{Alert Block}
        用于提示风险、限制、关键假设或需要听众立即注意的内容。
      \end{alertblock}

    \column{.31\textwidth}
      \begin{exampleblock}{Example Block}
        用于展示案例、经验规则、可复用写法或正向结论。
      \end{exampleblock}
  \end{columns}

  \vspace{.38cm}
  \begin{columns}[T,onlytextwidth]
    \column{.48\textwidth}
      \textbf{有序列表}
      \begin{enumerate}
        \item 第一层结论。
        \item 第二个步骤包含补充说明。
      \end{enumerate}
    \column{.48\textwidth}
      \textbf{无序列表与子项}
      \begin{itemize}
        \item 主题发现
        \item 证据组织
          \begin{itemize}
            \item 数据来源
            \item 分析方法
          \end{itemize}
      \end{itemize}
  \end{columns}
\end{frame}

\begin{frame}{定理、定义与证明环境}{沿用 Beamer 语义并接入主题配色}
  \begin{columns}[T,onlytextwidth]
    \column{.48\textwidth}
      \begin{nslrule}[证据链]
        若一个结论能够回溯到数据来源、处理步骤、指标定义和可视化解释，则称其具备可复核的生态证据链。
      \end{nslrule}

      \begin{nslclaim}[导航可解释性]
        当章节、小节和页面编号同时可见时，听众更容易在长报告中保持当前位置感。
      \end{nslclaim}

    \column{.48\textwidth}
      \begin{proof}
        封面后的目录提供全局结构，章节页提供局部结构，正文顶栏给出当前小节，底栏给出章节进度和页码。四类线索互相补充，因此可降低跳页、问答和回看时的定位成本。
      \end{proof}
  \end{columns}
\end{frame}

\subsection{公式表格}

\begin{frame}{公式与强调文本}{适合方法页和指标定义页}
  \scriptsize
  \begin{columns}[T,onlytextwidth]
    \column{.66\textwidth}
      \begin{block}{HP与CP的计算方法}
        $\romvar{Base}$为种族值，$\romvar{DV}$为个体值，$\romvar{STATEXP}$为基础点数，$\romvar{Level}$为等级，$\romvar{Stats}$为能力值，$\romvar{GoPower}$为GO力量。
      \end{block}
      \vspace{-.18cm}
      \begingroup
      \setlength{\abovedisplayskip}{.25em}
      \setlength{\belowdisplayskip}{.1em}
      \fontsize{6.15}{6.85}\selectfont
      \[
        \begin{aligned}
          \romvar{HP}
          &=\left\lfloor
          \frac{\left(\left(\romvar{Base}+\romvar{DV}\right)\times 2+
          \left\lfloor \frac{\sqrt{\romvar{STATEXP}}}{4} \right\rfloor\right)
          \times \romvar{Level}}{100}
          \right\rfloor+\romvar{Level}+10 \\[.28em]
          \romvar{CP}
          &=\min\Bigg(
          \Bigl\lfloor
          \left(\sum{\romvar{Stats}}-\sum{\romvar{GoPower}}\right)
          \times \frac{\romvar{Level}\times 6}{100} \\
          &\hspace{5.2em}
          +\sum{\romvar{GoPower}}\times
          \left(\frac{\romvar{Level}\times 4}{100}+2\right)
          \Bigr\rfloor,\;10000\Bigg)
        \end{aligned}
      \]
      \endgroup

    \column{.30\textwidth}
      \begin{exampleblock}{讲述方式}
        公式页不需要塞满推导；保留一个 \alert{可解释的核心式子}，再用简短文字说明变量含义。
      \end{exampleblock}
      \vspace{-.08cm}
      \begingroup
      \setbeamerfont{cookie quote}{size=\fontsize{7.3}{8.6},shape=\itshape}
      \setbeamerfont{cookie quote author}{size=\fontsize{5.8}{6.6},series=\bfseries}
      \begin{cookiequote}[模板提示]
        结构负责定位，内容负责说服，留白负责让听众呼吸。
      \end{cookiequote}
      \endgroup
  \end{columns}
\end{frame}

\begin{frame}{表格与宽版区域}{适合比较数据、实验设置和模板选项}
  \begin{wide}
    \begin{table}
      \centering
      \begin{tabularx}{.91\paperwidth}{lXXX}
        \toprule
        展示对象 & 推荐组件 & 适用场景 & 本模板效果 \\
        \midrule
        会议与机构信息 & 封面、顶栏、结尾页 & 学术报告、答辩、组会 & 统一露出会议名与中国科学院文献情报中心标识 \\
        章节结构 & 目录、章节页、小节页、底栏 & 多章节报告 & 自动显示当前位置和章节进度 \\
        关键结论 & block、alertblock、exampleblock & 结果页、方法页、结论页 & 与主色协调，避免额外自定义样式 \\
        代码和算法 & cookiecode、algorithm2e & 技术路线、复现说明 & 无需 shell escape，可直接编译 \\
        \bottomrule
      \end{tabularx}
    \end{table}
  \end{wide}
\end{frame}

\section{复杂版式}
\subsection{图示卡片}

\begin{frame}{TikZ 图示与卡片组合}{用于方法框架、流程图和模块关系}
  \begin{columns}[T,onlytextwidth]
    \column{.58\textwidth}
      \centering
      \begin{tikzpicture}[
        node distance=.62cm,
        every node/.style={font=\small},
        stage/.style={cookie accent node,minimum width=2.05cm,minimum height=.82cm,align=center},
        plain/.style={cookie node,minimum width=1.95cm,minimum height=.78cm,align=center},
        arrow/.style={cookie arrow}
      ]
        \node[stage] (a) {数据采集\\文献/专利/项目};
        \node[plain,right=of a] (b) {语义建模\\主题嵌入};
        \node[plain,right=of b] (c) {关系分析\\证据链};
        \node[stage,below=of b] (d) {情报呈现\\报告/图谱/建议};
        \draw[arrow] (a) -- (b);
        \draw[arrow] (b) -- (c);
        \draw[arrow] (b) -- (d);
        \draw[arrow] (c) -- (d);
      \end{tikzpicture}

    \column{.38\textwidth}
      \begin{cookiecard}[title=卡片一]
        卡片适合放置旁注、局部结论、数据口径或读图提示。
      \end{cookiecard}
      \begin{cookiecard}[title=卡片二]
        主题已内置节点、强调节点和箭头样式，适合快速绘制流程图。
      \end{cookiecard}
  \end{columns}
\end{frame}

\begin{frame}[
  bgimage=assets/文献中心logo.png,
  bgfit=contain,
  bgopacity=.10,
  bgoverlay=white,
  bgoverlayopacity=.62
]{背景图帧与局部视觉强调}{使用 frame key 控制普通页面背景}
  \begin{columns}[T,onlytextwidth]
    \column{.52\textwidth}
      \begin{block}{背景图参数}
        普通正文页可使用 \texttt{bgimage}、\texttt{bgfit}、\texttt{bgopacity}、\texttt{bgoverlay} 和 \texttt{bgoverlayopacity} 设置局部背景。
      \end{block}

    \column{.42\textwidth}
      \begin{exampleblock}{推荐用法}
        背景图只承担氛围和品牌提示，不应压过正文内容。正式报告中建议保持低透明度。
      \end{exampleblock}
  \end{columns}
\end{frame}

\subsection{代码动画}

\begin{frame}[fragile]{代码框与配置片段}{cookiecode 基于 listings 与 tcolorbox}
  \begin{cookiecode}[title=template-config.py,listing options={style=cookiebare,language=Python,basicstyle=\ttfamily\scriptsize}]{Python}
theme = "cookie"
accent = "#642860"
conference = "这里放会议名称/硕博论文答辩"
logo = "assets/nsl-logo.png"
footLogo = "assets/nsl-logo-white.png"

sections = [
    "模板全局结构",
    "内容组件",
    "复杂版式",
    "引用与收束",
]
  \end{cookiecode}
\end{frame}

\begin{frame}{叠加动画与模态层}{覆盖层用于逐步讲解复杂关系}
  \begin{columns}[T,onlytextwidth]
    \column{.54\textwidth}
      \begin{enumerate}
        \item \cookiefocus<1>{第一步：确定问题边界。}
        \item \cookiefocus<2>{第二步：补充数据与指标。}
        \item \cookiefocus<3>{第三步：聚焦结论与解释。}
      \end{enumerate}
      \cookiereveal<2->{%
        \begin{block}{逐步显示}
          \texttt{\textbackslash cookiereveal} 会保留版面空间，适合在同一页逐步补充内容。
        \end{block}
      }

    \column{.42\textwidth}
      \begin{cookiecard}[title=覆盖层状态]
        本页共有三个覆盖层。导出 PDF 后，动画会展开为连续页面，便于演示和打印。
      \end{cookiecard}
  \end{columns}

  \begin{modal}<3>[模态层示例]
    当前覆盖层可以把注意力集中到一个关键解释上，例如“为什么这一页需要暂停讨论”。
  \end{modal}
\end{frame}

\begin{frame}{算法环境}{基于强化学习的双打 OU 对战智能体训练算法}
  \vspace{.04cm}
  \begin{wide}
    \begin{algorithm}[H]
      \caption{基于强化学习的双打 OU 对战智能体训练算法}
      \KwIn{初始策略网络 $\pi_\theta$，价值网络 $V_\phi$，训练环境 $E$，经验回放池 $D$，最大训练轮数 $N$，折扣因子 $\gamma$，学习率 $\alpha$}
      \KwOut{训练完成的双打 OU 对战智能体 $\pi_\theta^*$}
      初始化策略网络参数 $\theta$ 与价值网络参数 $\phi$\;
      初始化经验回放池 $D$\;
      构建双打 OU 对战环境 $E$，包括队伍信息、招式信息、属性克制关系与场地状态\;
      \For{训练轮数 $i=1$ \KwTo $N$}{
        重置环境 $E$，获得初始状态 $s_0$\;
        \While{当前对战未结束}{
          根据当前状态 $s_t$，由策略网络 $\pi_\theta$ 生成动作概率分布\;
          从合法动作集合中选择动作 $a_t$\;
          在环境中执行动作 $a_t$，获得下一状态 $s_{t+1}$、奖励 $r_t$ 与终止标记 $done$\;
          将转移样本 $(s_t,a_t,r_t,s_{t+1},done)$ 存入经验回放池 $D$\;
          根据场上优势、剩余宝可梦数量、造成伤害、击倒对手和最终胜负计算综合奖励\;
          令 $s_t \leftarrow s_{t+1}$\;
        }
        从经验回放池 $D$ 中采样一批训练样本\;
        计算目标价值 $y_t=r_t+\gamma V_\phi(s_{t+1})$\;
        计算优势函数 $A_t=y_t-V_\phi(s_t)$\;
        根据优势函数更新策略网络参数 $\theta$，提高高收益动作的选择概率\;
        根据目标价值 $y_t$ 更新价值网络参数 $\phi$，降低状态价值预测误差\;
        定期将当前智能体与历史版本智能体进行自博弈，以提升策略稳定性\;
        定期在固定测试队伍和随机对手队伍上评估智能体胜率\;
      }
      \KwRet 训练后的策略网络 $\pi_\theta^*$\;
    \end{algorithm}
  \end{wide}
\end{frame}

\cookieplain{一句话强调页：把关键结论留给一整页。}

\section{引用与收束}
\subsection{参考文献}

\begin{frame}{引用与参考文献样式}{BibLaTeX 负责数据，主题负责视觉}
  \begin{columns}[T,onlytextwidth]
    \column{.52\textwidth}
      \begin{block}{引用示例}
        Beamer 的结构化页面适合技术报告和答辩场景 \parencite{beameruserguide}。本模板的视觉方向参考了 Cookie Beamer \parencite{cookiebeamer}，同时保留对 Metropolis 和 awesome-beamer 设计习惯的说明 \parencite{metropolis,awesomebeamer}。
      \end{block}

    \column{.42\textwidth}
      \begin{cookiecard}[title=参考文献页]
        文献数据写在 \texttt{refs.bib} 中；编译时运行 LuaLaTeX 与 Biber，即可生成文末参考文献。
      \end{cookiecard}
  \end{columns}
\end{frame}

\begin{frame}[allowframebreaks]{参考文献}
  \printbibliography[heading=none]
\end{frame}

\subsection{使用建议}

\begin{frame}{正式使用前的替换清单}{保留结构，替换内容}
  \begin{columns}[T,onlytextwidth]
    \column{.48\textwidth}
      \begin{block}{必须替换}
        \begin{enumerate}
          \item 论文题名、副标题、作者和单位。
          \item 会议名称或答辩场景名称。
          \item 日期、通讯作者邮箱和参考文献数据。
        \end{enumerate}
      \end{block}

    \column{.48\textwidth}
      \begin{exampleblock}{建议保留}
        \begin{enumerate}
          \item 章节和小节命令，用于驱动导航。
          \item 文献中心紫色主色与 Logo。
          \item block、卡片、表格、代码框等主题组件。
        \end{enumerate}
      \end{exampleblock}
  \end{columns}
\end{frame}

\end{document}
